\chapter{The enveloping bar and the audit of the planar tables}
\label{ch:enveloping-audit}

The three bars of the typed bar theorem are computed here on enveloping
algebras, and the computation settles two dimension formulas that circulated
as bar cohomology.

\section{Chevalley comparison}

\begin{theorem}[Bar--Chevalley comparison]
\label{thm:bar-chevalley-spine}
Let $\g$ be a Lie algebra object in $\FactD(X)$ for the pointwise product,
$K$-flat, positively graded and complete, with every total weight admitting
finitely many decompositions into positive weights, and suppose internal PBW
holds for $U(\g)$. There is a natural quasi-isomorphism of conilpotent
coalgebras in $\FactD(X)$
\[
 \BarPerp\bigl(U(\g)\bigr)\ \xrightarrow{\ \sim\ }\ C_*(\g),
\]
compatible with the factorization coefficient and the PBW filtration.
\end{theorem}
\begin{proof}
Filter $U(\g)$ by PBW degree and both complexes by total word length. The
associated graded enveloping algebra is $\Sym(\g)$, and the bar of a symmetric
algebra is the suspended symmetric coalgebra on $\g$, which is the associated
graded of Chevalley chains. The comparison is an isomorphism on the first page;
finite decomposition gives a bounded exhaustive spectral sequence in each
weight, and completeness promotes it to a quasi-isomorphism. Every step uses
only internal tensor products and the bracket, so the map stays in $\FactD(X)$.
\end{proof}

\begin{calculation}[Finite-dimensional simple case]
\label{calc:simple-spine}
For a simple Lie algebra of rank $r$ with exponents $m_1,\dots,m_r$, the
cohomology is the exterior algebra on primitive classes of degrees $2m_i+1$:
\[
 \BarPerp U(\mathfrak{sl}_2):\ 1+t^3,
 \qquad
 \BarPerp U(\mathfrak{sl}_3):\ 1+t^3+t^5+t^8,
 \qquad
 \BarPerp U(\mathfrak g_2):\ 1+t^3+t^{11}+t^{14}.
\]
The first two were recomputed by exact rational elimination on
$\Lambda^\bullet\g^*$ built from matrix units, with $d^2=0$ asserted on every
cochain space rather than assumed, giving Betti numbers $1,0,0,1$ and
$1,0,0,1,0,1,0,0,1$. The Chevalley chain dimensions are
$\binom{\dim\g}{n}$, so for $\mathfrak{sl}_2$ they are $1,3,3,1$ while the bar
chain dimensions of an augmented algebra with $\dim\bar A=3$ are $3^r$.
\end{calculation}

\section{The Virasoro row}

The transverse presentation of the Virasoro family is the enveloping algebra of
the positive Witt algebra
\[
 L_1=\bigoplus_{n\ge1}\mathbb C\,e_n,
 \qquad
 [e_i,e_j]=(j-i)e_{i+j},
\]
to which Theorem~\ref{thm:bar-chevalley-spine} applies: $L_1$ is positively
graded with one-dimensional weight spaces, so every weight has finitely many
decompositions.

\begin{theorem}[Pentagonal rigidity]
\label{thm:pentagonal-spine}
The bar homology of $U(L_1)$ satisfies
\[
 \dim H^n\bigl(\BarPerp U(L_1)\bigr)=2
 \qquad\text{for every }n\ge1,
\]
one dimension in each of the weights $(3n^2-n)/2$ and $(3n^2+n)/2$. Its graded
Euler character is
\begin{equation}
\label{eq:witt-euler-spine}
 \chi_q\bigl(\BarPerp U(L_1)\bigr)=\prod_{n\ge1}\bigl(1-q^n\bigr),
\end{equation}
and \eqref{eq:witt-euler-spine} is accounted for degree by degree with no
residual freedom.
\end{theorem}
\begin{proof}
By Theorem~\ref{thm:bar-chevalley-spine} the bar of $U(L_1)$ is $C_*(L_1)$,
whose cohomology is Goncharova's: two dimensional in each degree $n\ge1$, in
the weights $(3n^2\mp n)/2$. Identity \eqref{eq:witt-euler-spine} is the
denominator product applied to $L_1$, where every root space is one dimensional
and even, so $\smult(L_1)_n=1$. Euler's pentagonal number theorem expands the
product as $\sum_{m\in\Z}(-1)^mq^{m(3m-1)/2}$, whose support is exactly
$\{(3n^2\pm n)/2\}$ with coefficient $(-1)^n$ at each of $(3n^2-n)/2$ and
$(3n^2+n)/2$. Comparing with $\sum_n(-1)^n\dim H^n_w$, each weight of the Euler
character is saturated by a single class in the degree Goncharova predicts, so
no further classes and no cancellation are available.
\end{proof}

The cohomology was computed for $n=1,2,3,4$, occupying the weights $(1,2)$,
$(5,7)$, $(12,15)$, $(22,26)$, one dimension in each; the nonzero coefficients
of $\prod_{n\ge1}(1-q^n)$ occupy $0,1,2,5,7,12,15,22,26$ with signs
$+,-,-,+,+,-,-,+,+$.

\section{The planar tables}

Two formulas circulated as bar cohomology dimensions:
$\dim H^n=M(n+1)-M(n)=1,2,5,12,30,76,\dots$ for the Virasoro family, with $M$
the Motzkin numbers, and $\dim H^n=R(n+3)=3,6,15,36,91,\dots$ for the
$\mathfrak{sl}_2$ family, with $R$ the Riordan numbers. The coincidence of the
discriminant $\sqrt{1-2z-3z^2}$ in the two generating functions was attributed
to Drinfeld--Sokolov reduction.

\begin{proposition}[The two sequences are one sequence]
\label{prop:one-sequence-spine}
For all $n\ge0$,
\[
 M(n)=R(n)+R(n+1),
 \qquad
 M(n+1)-M(n)=R(n+2)-R(n).
\]
The Virasoro sequence is therefore $R(n+2)-R(n)$ and the $\mathfrak{sl}_2$
sequence is $R(n+3)$: both are elementary transforms of one sequence, and their
generating functions are rational in the single algebraic function
$\sqrt{1-2z-3z^2}$ for that reason alone. No Drinfeld--Sokolov input is used or
needed.
\end{proposition}
\begin{proof}
Riordan numbers count Motzkin paths with no level step at height zero.
Splitting a Motzkin path of length $n$ according to whether its first step is
level gives $M(n)=R(n)+R(n+1)$; subtracting consecutive instances gives the
second identity. Both were checked termwise for $n<24$.
\end{proof}

\begin{theorem}[Carrier theorem for the planar series]
\label{thm:carrier-spine}
Neither sequence is the ordered bar homology of the algebra it names, and
neither is a bar chain count.
\begin{enumerate}
\item For $\mathfrak{sl}_2$ the bar homology is $1,0,0,1$ by
  Calculation~\ref{calc:simple-spine}, and the bar chain dimensions are
  $3,9,27,81,243$. The sequence $3,6,15,36,91$ is neither.
\item For the Virasoro family the bar homology is the constant $2$ by
  Theorem~\ref{thm:pentagonal-spine}. On degrees $n=1,2,3$ the asserted values
  are $1,2,5$; they agree with the truth at $n=2$ and nowhere else.
\item Independently of any presentation, a Virasoro chiral coefficient does not
  determine the transverse bar: a fixed coefficient may be combined with free,
  square-zero, enveloping, or quantum-plane transverse algebras whose bars
  differ, while the coefficient and the central charge are unchanged.
\end{enumerate}
Consequently a chain-level identification of either planar series with an
ordered bar requires a specified transverse multiplication together with a
quasi-isomorphism from its planar complex to that bar, and by (i) and (ii) no
such quasi-isomorphism exists for the algebras named.
\end{theorem}

\begin{remark}[The type error beneath the numerical one]
The bar homology of $U(L_1)$ is bigraded by cohomological degree and weight,
and is finite in each degree only because Theorem~\ref{thm:pentagonal-spine}
concentrates it in two weights. A formula assigning one integer to each degree,
with no weight truncation named, does not specify an invariant of a graded
object with infinite-dimensional graded pieces.
\end{remark}

\chapter{The obstruction constitution}
\label{ch:obstructions}

Each statement below names what is \emph{not} determined, and therefore what
must be supplied. Without them the preceding theorems admit false corollaries.

\begin{proposition}[Where a mixed distributive law can live]
\label{prop:lambda-sector-spine}
The sector on which a mixed distributive law between $\tensor^!$ and $\starCh$
is geometrically defined is the common-decomposition sector, the image of the
aligned idempotent $e_{\mathrm{al}}$. Off that sector the two products admit no
interchange, only the aligned retraction, so no distributive law can be
written.
\end{proposition}
\begin{proof}
A distributive law requires a natural transformation interchanging the two
products in a fixed direction, subject to the coherence conditions of each. The
canonical comparison in the additive model is the aligned inclusion $\iota$ with
retraction $\pi$, and $\pi\iota=\id$ while $\iota\pi=e_{\mathrm{al}}\ne\id$ in
general. A natural transformation subject to both coherence conditions
therefore exists only where $e_{\mathrm{al}}$ acts as the identity.
\end{proof}

This locates a mixed law without constructing one. No example in the present
corpus exhibits one; Zamolodchikov--Faddeev algebras are the natural candidate,
since their spectral variable belongs to the chiral direction while their word
belongs to the transverse direction.

\begin{theorem}[No quantum group from dimension loss]
\label{nogo:dimension-loss-spine}
The dimension of the kernel of $V^{\otimes n}\to\Sym^nV$ determines no
$R$-matrix, associator, or quantum group. That kernel carries no
multiplication, coproduct, monodromy, or coherence equation.
\end{theorem}
\begin{proof}
The kernel is a representation of $\Sigma_n$ and nothing more. Two algebras
with the same underlying graded vector space, hence the same kernel dimensions
in every arity, may have different multiplications and inequivalent
representation categories; a free algebra and a quadratic quotient with the
same Hilbert series is an example. A quantum-group structure is a coproduct
together with coherence data, none of which is a function of a dimension.
\end{proof}

\begin{theorem}[No modular object from topology alone]
The identity $\partial^2=0$ on $\overline{\mathcal M}_{g,n}$ assigns no
coefficient operations to boundary strata. A modular master element exists only
after cyclic contractions and clutching-compatible coefficient maps have been
supplied.
\end{theorem}

\begin{theorem}[No bulk from a boundary centre]
The derived centre is terminal among algebraic central actions on the boundary
algebra. It does not detect a decoupled bulk sector, so a physical bulk is not
determined by the boundary chiral algebra unless a bulk-to-centre comparison is
supplied and proved to be an equivalence.
\end{theorem}
\begin{proof}
Terminality is the universal property of the centre. It is a statement about
the boundary category and is insensitive to any sector acting trivially on it,
so a bulk with trivial boundary action is invisible.
\end{proof}

\begin{theorem}[No classification by central charge]
Central charge determines the Virasoro anomaly coefficient. It does not
determine the chiral operation, the representation category, the braid
monodromy, or the higher-genus conformal blocks. Distinct holomorphic theories
at equal central charge exist.
\end{theorem}

\begin{theorem}[No identification of anomalies of different types]
\label{nogo:anomalies-spine}
A current-algebra level, a Virasoro central charge, a determinant-line
curvature, a BRST ghost-number anomaly, and a Borcherds-product weight are
invariants in different theories, living in different deformation complexes.
They may be related by a theorem in a particular family. They are not one
universal scalar, and no identity among them follows from numerical agreement.
\end{theorem}
\begin{proof}
Each is a class in a distinct complex: an Atiyah-algebra extension class, a
Virasoro cocycle, the curvature of a determinant line, a class in $H^1$ or
$H^2$ of a BRST convolution complex, and a weight in the character lattice of
an automorphic form. A map between two of these complexes is additional data;
in its absence, equality of numbers is a coincidence of values, not an
identification of classes.
\end{proof}

\begin{theorem}[No unnamed comparison of bars]
An asserted equality $\BarPerp(A)=\BarAssCh(A)$ has no meaning without a
functor carrying the pointwise multiplication problem to the chiral-convolution
multiplication problem. The two have different underlying gradings, counting
interval subdivisions and collision trees respectively. Agreement of Euler
characteristics is not a comparison map.
\end{theorem}

\begin{theorem}[A scalar sequence is not homology data]
\label{nogo:scalar-spine}
A sequence of integers becomes bar homology only after a chain complex, a
differential, a grading convention, and a rank computation or quasi-isomorphism
are specified. Matching initial values, discriminants, or growth rates
constructs no such map.
\end{theorem}
\begin{proof}
Chain-space dimensions and homology dimensions differ by the ranks of the
differentials, which are not functions of the chain dimensions.
Calculation~\ref{calc:simple-spine} exhibits an algebra whose bar chain
dimensions are $3,9,27,81,243$ and whose bar homology dimensions are $1,0,0,1$;
Proposition~\ref{prop:one-sequence-spine} exhibits two sequences with the same
discriminant and growth rate that are transforms of one another and of neither
homology.
\end{proof}

Theorem~\ref{nogo:anomalies-spine} together with the central-charge statement
retracts any programme treating a tuple of scalars drawn from these sources as
a single classifying invariant, and any trace identity asserted among them
without a named comparison morphism.
